% arXiv-ready LaTeX (cs.MA target)
\documentclass[11pt]{article}
\usepackage[utf8]{inputenc}
\usepackage[T1]{fontenc}
\usepackage{times}
\usepackage{fullpage}
\usepackage{graphicx}
\usepackage{booktabs}
\usepackage{amsmath,amssymb}
\usepackage{hyperref}
\usepackage{enumitem}
\usepackage{microtype}
\usepackage{siunitx}
\usepackage{caption}
\usepackage{float}

\title{BVR-Sim: Advancing Open Beyond-Visual-Range Air Combat Environments for Reinforcement Learning}

\author{Haocheng Sun\thanks{sunhaocheng@bupt.edu.cn} \\ Beijing University of Posts and Telecommunications}
\date{}

\begin{document}
\maketitle

\begin{abstract}
Reinforcement learning (RL) for beyond-visual-range (BVR) air combat requires environments that expose realistic sensing, guidance, and control constraints while remaining open and reproducible. Existing open simulators often simplify missile guidance, omit radar/datalink modeling, or lack the instrumentation needed for large-scale RL workflows. We present \textbf{BVR-Sim}, a fully open environment integrated into a multi-team RL stack that unifies high-fidelity aircraft and missile dynamics, configurable sensor/datalink modeling, autopilot-aware action abstractions, and comprehensive training instrumentation including trajectory logging and reward visualization. BVR-Sim aims to lower the barrier to reproducible BVR research by providing an end-to-end pipeline from scenario configuration to ACMI/trajectory export. We describe the system architecture, physics and sensing models (including Mach-indexed drag and hybrid guidance), observation/action interfaces, and practical instrumentation for training and offline data collection. We include placeholders for experimental protocols and evaluation metrics that are suitable for benchmarking self-play, distillation, and offline RL on BVR tasks.
\end{abstract}

\section{Introduction}
Beyond-visual-range (BVR) air combat couples multi-agent tactics, imperfect sensing, and constrained missile guidance envelopes. Studying RL approaches in this domain requires environments that faithfully expose the dynamics and sensing constraints that shape decision-making while remaining open and instrumented for reproducible research. Many released environments fall into two camps: (1) high-fidelity but closed military simulators, and (2) lightweight open frameworks that simplify aerodynamics and guidance (often to proportional navigation). The former are inaccessible to most researchers; the latter omit critical phenomena such as Mach-dependent missile drag, datalink noise, and energy-aware launch trade-offs.

We present \textbf{BVR-Sim}, an open system designed as a research-grade infrastructure for RL and imitation learning in BVR engagements. BVR-Sim emphasizes (i) physically motivated missile and aircraft solvers, (ii) flexible sensor/datalink abstractions, (iii) autopilot-aware action interfaces so agents reason at a tactical level, and (iv) full instrumentation for debugging, benchmarking, and offline data export. Our contributions are: a system design and implementation that integrates these components into a single training loop; exposed, testable physics and guidance models (including Mach-indexed drag and hybrid guidance blending angle and LOS-rate control); and a suite of training hooks (reward visualizer, ACMI export, trajectory recorder) to accelerate reproducible research.

\section{Related work}
Open RL environments for aerial combat span a range of fidelity and usability trade-offs. \textbf{BVR-Gym} employs JSBSim for aircraft dynamics and simplifies missile guidance to proportional navigation (PN), focusing on accessible aircraft controllers but omitting advanced sensing and instrumentation. \textbf{B-ACE} prioritizes accessibility via Godot but trades off aerodynamic realism and missile detail. \textbf{WUKONG} demonstrates RL-driven tactics at scale but is not fully open and lacks public instrumentation for replay and debug. Our work sits between these extremes: an open, instrumented platform that preserves physically meaningful missile and sensing models while remaining practical for RL workflows.

\section{System overview and design goals}
BVR-Sim is designed with three guiding principles:
\begin{enumerate}[noitemsep]
    \item \textbf{Physical fidelity where it matters.} Represent missile aerodynamics (Mach-dependent drag, gravity), energy-aware launch conditions, and guidance laws that influence tactical outcomes.
    \item \textbf{Action abstraction for learning.} Provide autopilot-aware, interpretable action spaces so agents learn tactics rather than low-level control.
    \item \textbf{Instrumentation and reproducibility.} Built-in logging, ACMI export, reward decomposition, and trajectory recording for offline RL and reproducible benchmarks.
\end{enumerate}

At a high level the mission module (e.g., \texttt{MISSION/bvr\_sim\_v3}) exposes a Gym-like environment (here called \texttt{BVR3DEnv}) wrapped by a multi-team runner. The configuration is JSONC-driven to allow scenario variants without code changes (spawn geometry, observation presets, reward weights, logging paths).

% % FIGURE: system architecture
% \begin{figure}[H]
%   \centering
%   \fbox{\begin{minipage}[c][5cm][c]{0.9\linewidth}
%   \centering \textbf{Figure placeholder: System architecture diagram (mission runner, env wrapper, physics, sensing, recorder, visualizer).}
%   \end{minipage}}
%   \caption{System architecture (placeholder).}
%   \label{fig:architecture}
% \end{figure}

\section{Environment interface}
\subsection{Action and control abstraction}
Agents act in discrete MultiDiscrete spaces that map to autopilot-level rate commands. A typical action at time $t$ is
\[
\mathbf{a}_t = (k_\psi, k_h, k_v, k_\ell),
\]
where discrete indices $k_\psi,k_h,k_v$ encode heading, altitude and speed rate bins, and $k_\ell$ is a binary launch command. The mapping to continuous rate deltas is:
\begin{align}
\Delta \psi_t &= \gamma_\psi \frac{k_\psi-\bar{k}_\psi}{\bar{k}_\psi}, &
\Delta h_t &= \gamma_h \frac{k_h-\bar{k}_h}{\bar{k}_h}, &
\Delta v_t &= \gamma_v \frac{k_v-\bar{k}_v}{\bar{k}_v}, \label{eq:action_map}
\end{align}
with midpoints $\bar{k}_\psi$, $\bar{k}_h$, $\bar{k}_v$ and scaling constants $\gamma_{\{\psi,h,v\}}$ (configurable via JSONC). The launch command is gated by radar lock and other safety checks.

\subsection{Observation encoding}
Observation spaces are configurable (compact, extended, shadow) and provide normalized relative kinematics for self, allies, enemies, and up to a fixed number of in-flight friendly missiles. Observations include noisy positional/velocity estimates, missile-warning booleans, lock indicators, and time-to-impact estimates. A concatenated observation can be written as:
\begin{align}
o_t &= \big[s^{\text{self}}_t,\ \{s^{\text{ally}}_{t,i}\}_{i=1}^{N_a},\ \{s^{\text{enemy}}_{t,j}\}_{j=1}^{N_e},\ \{s^{\text{miss}}_{t,m}\}_{m=1}^{M}\big]. \label{eq:obs_concat}
\end{align}
Sensing noise is modeled as additive Gaussian perturbations on position and velocity:
\begin{align}
\tilde{\mathbf{p}}_{t} &= \mathbf{p}_{t} + \boldsymbol{\epsilon}_p,&
\tilde{\mathbf{v}}_{t} &= \mathbf{v}_{t} + \boldsymbol{\epsilon}_v, \\
\boldsymbol{\epsilon}_p &\sim \mathcal{N}(0,\sigma_p^2 I_3),\quad
\boldsymbol{\epsilon}_v \sim \mathcal{N}(0,\sigma_v^2 I_3). \label{eq:sensing_noise}
\end{align}

Episode termination enforces that all in-flight missiles have resolved before ranking teams to avoid bias against long time-of-flight intercepts:
\begin{align}
\text{done}_t &= \mathbb{1}\!\big((n_{\text{red},t}=0 \lor n_{\text{blue},t}=0)\land|\mathcal{M}_t|=0\big)\ \lor\ (t\ge T_{\max}). \label{eq:termination}
\end{align}

\section{Physics and guidance models}
\subsection{Aircraft and missile dynamics}
Aircraft evolve through a pluggable flight-dynamics model:
\begin{align}
\mathbf{x}_{t+1} &= \mathbf{x}_t + f_{\text{FDM}}\big(\mathbf{x}_t, \mathbf{u}_t(\Delta\psi_t,\Delta h_t,\Delta v_t)\big)\Delta t, \label{eq:fdm}
\end{align}
where state $\mathbf{x}$ aggregates position, velocity and attitude, and $\mathbf{u}$ is the autopilot command derived from the rate deltas in Eq.~\eqref{eq:action_map}. The modularity of $f_{\text{FDM}}$ allows swapping flight-dynamics models without changing the RL interface.

Missile dynamics include thrust, Mach-dependent drag, gravity, and aerodynamic losses, integrated synchronously with aircraft solvers so that time alignment with the RL loop is preserved. Drag is computed as:
\begin{align}
D &= \tfrac{1}{2}\rho(h)v^2 S_{\text{ref}} C_x(M,\alpha), \label{eq:drag} \\
C_x(M,\alpha) &= C_{x0}(M) + C_{xB}(M) + K_1(M)|\alpha| + K_2(M)\alpha^2 + C_{x,\text{wave}}(M) + C_{x,\text{sup}}(M). \label{eq:cx}
\end{align}
Here $M=v/a(h)$ is Mach number, $\alpha$ the angle-of-attack, $\rho(h)$ altitude density, and lookup tables supply $C_{x0},C_{xB},K_1,K_2$. A transonic wave-drag bump and supersonic correction are included to match published performance envelopes where available.

\subsection{Guidance}
Guidance is modeled using commanded load factors derived from a blend of angle-error and LOS-rate controllers. The horizontal LOS-rate command is:
\begin{align}
\dot{\beta}_{\text{cmd}} &= \lambda(t)\,k_\beta\big[\phi - (\beta - \pi)\big] + (1-\lambda(t))\,\dot{\beta}_{\text{meas}}, \label{eq:beta_cmd}
\end{align}
where $\lambda(t)=\max((t_{\text{thrust}}-t)/t_{\text{thrust}},0)$ blends angle-error steering during motor burn into rate-based steering thereafter. The commanded lateral/vertical load factors are:
\begin{align}
n_y &= \frac{K(R)}{g}\,v\cos\theta\,\dot{\beta}_{\text{cmd}}, \qquad
n_z = \frac{K(R)}{g}\,v\,\dot{\epsilon} + \cos\theta, \label{eq:load}
\end{align}
with gain $K(R)$ range-dependent (ramped up when the radar sees the target within a threshold). Load factors saturate at platform limits (e.g., $\pm 10$ g), and the exposed formulas let researchers replace the guidance law (e.g., augmented PN, differential-game controllers) while holding sensing and drag constant for controlled ablations.

\section{Training instrumentation and workflows}
\subsection{Reward composition and visualization}
Rewards are decomposed into dense shaping and sparse event terms:
\begin{align}
R_t &= \sum_{c\in\mathcal{C}} w_c r^c_t, \qquad \mathcal{C}=\{\text{engage},\text{altitude},\text{safe-alt},\text{launch},\text{hit},\text{win},\ldots\}, \label{eq:reward_sum}
\end{align}
with weights configured via JSONC. A per-component visualizer emits JSON and PNG plots (symlog scale) each episode, enabling diagnosis of mode collapse or overly sparse signals.

\subsection{Distillation and teacher forcing}
BVR-Sim supports distillation and teacher-forcing hooks where a scripted tactical baseline can provide action targets:
\begin{align}
\mathcal{L}_{\text{dist},t} &= \lVert \hat{\mathbf{a}}^{\text{RL}}_t - \hat{\mathbf{a}}^{\star}_t\rVert_p + w_{\text{shoot}}|\sigma^{\text{RL}}_t - \sigma^{\star}_t|, \label{eq:distill_loss}\\
\tilde{\mathbf{a}}_t &= (1-\eta_t)\mathbf{a}^{\text{RL}}_t + \eta_t\mathbf{a}^{\star}_t, \quad \eta_t\in\{0,1\}, \label{eq:teacher_force}
\end{align}
where $\mathbf{a}^\star$ is a baseline action (scripted policy) and $\eta_t$ toggles teacher forcing when enabled. This supports curricula that anneal imitation weight while preserving robust fallback behavior in training.

\subsection{Baselines and fallback control}
Scripted opponents (random, simple heuristics, tactical state machine) are provided and run inside the same mission process to enable comparable instrumentation and debugging. Baseline policies expose their internal decision signals to the reward visualizer, facilitating comparisons with learned policies.

\subsection{Trajectory recording for offline RL}
Aircraft and missile-level traces are recorded into chunked replay pools suitable for offline RL and behavior cloning. Each record entry contains:
\[
(o_t, \mathbf{a}_t, r_t, o_{t+1}, \text{aux}_t),
\]
where auxiliary info includes Mach number, autopilot deltas, shoot flags, and physics metadata. The recorder supports episodic flushing and safe archival to make the dataset immediately usable for downstream offline experiments.

% % TABLE: feature comparison (placeholder)
% \begin{table}[H]
%   \centering
%   \caption{Feature comparison across open BVR environments (placeholder).}
%   \label{tab:comparison}
%   \fbox{\begin{minipage}{0.95\linewidth}
%   \centering Table placeholder: Fill with quantitative feature matrix (e.g., dynamics fidelity, guidance detail, sensing, datalink, logging). 
%   \end{minipage}}
% \end{table}

\section{Evaluation (protocols and metrics)}
This work is primarily a systems/infrastructure contribution; the open environment is intended to support a range of experimental studies. Below we specify recommended benchmarking protocols and metrics that can be used to evaluate algorithms on BVR-Sim.

\paragraph{Suggested scenarios} Provide reproducible scenario configurations (via JSONC) with varying separation, altitude/heading spread, and force composition (1v1, 2v2). Default spawn randomization uses \SI{37.2}{NM} mean separation with \SI{2}{NM} spread.

\paragraph{Metrics} Primary metrics include team win rate, average time-to-engagement, missile survival rate, and energy-state statistics (e.g., post-encounter speed/altitude distributions). For learning algorithms measure sample efficiency (episodes-to-threshold win rate), robustness across seeds, and generalization to different team sizes.

\paragraph{Baselines and experiments} We recommend evaluating:
\begin{itemize}[noitemsep]
  \item Scripted baselines (heuristic state machine) as an absolute floor.
  \item PPO/A2C-style on-policy learners with the autopilot-aware action space.
  \item Distillation experiments where imitation weight is annealed.
  \item Offline RL / behavior cloning using recorded trajectory pools.
\end{itemize}

% % FIGURE: sample learning curves placeholder
% \begin{figure}[H]
%   \centering
%   \fbox{\begin{minipage}[c][4cm][c]{0.9\linewidth}
%   \centering \textbf{Figure placeholder: Example learning curves (win rate, reward components) and profiler stage timings.}
%   \end{minipage}}
%   \caption{Example evaluation artifacts (placeholders).}
%   \label{fig:eval}
% \end{figure}

\section{Implementation notes and reproducibility}
The system is driven by a JSONC configuration system that declares spawn geometry, opponent selection, observation presets, reward weights, and logging paths. Example commands to run training:
\begin{verbatim}
python main.py --cfg MISSION/bvr_sim_v3/conf_system/<name>.jsonc
\end{verbatim}
Run outputs auto-create ACMI dumps, per-component reward plots (JSON + PNG), and trajectory archives under the configured log directory. Timing and profiling hooks are available (e.g., \texttt{StepProfiler}) and emit running means and standard deviations for pipeline stages:
\begin{align}
\mu_k &= \frac{1}{N_k}\sum_{i=1}^{N_k}\Delta t_{k,i}, &
\sigma_k &= \sqrt{\frac{1}{N_k-1}\sum_{i=1}^{N_k}(\Delta t_{k,i}-\mu_k)^2}. \label{eq:profiler}
\end{align}

\paragraph{Code references vs descriptive paths} To improve readability in the paper we use descriptive path references (e.g., "mission module directory", "simulator/sense") but retain example paths in parentheses for reproducibility. If you would like the manuscript to include raw code paths verbatim everywhere, I can reintroduce them.

\section{Discussion and limitations}
BVR-Sim emphasizes system-level fidelity and instrumentation; however, it is not a drop-in replacement for classified, high-fidelity defense simulators. Limitations include the finite set of missile/aircraft models provided and assumptions in sensing/datalink noise models. Future work includes expanding the model library, adding formation-level behaviors, and publishing benchmark splits and recorded datasets for community comparison.

\section{Conclusion}
We introduced BVR-Sim, an open, instrumented environment for BVR air combat research that combines higher-fidelity missile aerodynamics, configurable sensing/datalink models, autopilot-aware actions, and a complete training/recording toolchain. By providing a reproducible, extensible infrastructure, BVR-Sim aims to make realistic BVR RL research more accessible to the community.

% \section*{Acknowledgements}
% (Placeholder) — acknowledge collaborators, funding, and code contributors here.

\begin{thebibliography}{9}
\bibitem{scukins2024bvrgym}
E.~Scukins, M.~Klein, L.~Kroon, and P.~\"Ogren, ``BVR Gym: A Reinforcement Learning Environment for Beyond-Visual-Range Air Combat,'' \emph{arXiv preprint} arXiv:2403.17533, 2024.

\bibitem{kuroswiski2024bace}
A.~R. Kuroswiski, A.~S. Wu, and A.~Passaro, ``B-ACE: An Open Lightweight Beyond Visual Range Air Combat Simulation Environment for Multi-Agent Reinforcement Learning,'' 2024. DOI: 10.13140/RG.2.2.11999.57762.

\bibitem{piao2020wukong}
H.~Piao \emph{et~al.}, ``WUKONG: Beyond-Visual-Range Air Combat Tactics Auto-Generation by Reinforcement Learning,'' \emph{IEEE}, 2020.
\end{thebibliography}

% % Appendices and extended derivations (optional)
% \appendix
% \section{Appendix: Example configuration keys}
% (Placeholder) Provide examples of JSONC keys, observation presets, and reward-weight defaults here.

\end{document}
